\section{That the chain is not vacuous, and what it still does not say}
\label{sec:satisfiability}
Both chain theorems are conditional: each accepts eight named implications and
concludes \texttt{UnifiedSelf}. A conditional theorem whose hypotheses nothing
satisfies proves nothing about anything, so the development supplies witnesses
discharging all eight at once, and --- separately --- counterexamples showing
that the same hypotheses cannot be met by substituting an arbitrary target.

\subsection{All eight arrows on one substrate}
\begin{claim}
There are a double well, an absorbing register, an erasing update, a refining
mesh, and a field and cover over the type \texttt{Cortex} that satisfy $E12$
through $E89$ simultaneously; hence \texttt{UnifiedSelf} holds for the supplied
reflexive boundary.
\end{claim}
\begin{proof}
Apply the passive chain to those components. The first four arrows are supplied
by their witnesses; $E56$ is discharged by the numerals $D=1$, $K=3$, $L=3$
together with the identification of the supplied field, $E67$ by evaluating the
critical coupling and comparing $3>2$, $E78$ by the cover's reachability at
coupling three, and $E89$ by the constructed section, its restriction equations,
the Lipschitz bound and the fixed-point equation.
\end{proof}
\formal{chain-hypotheses-jointly-satisfiable}
\scope The components are separate toy constructions sharing one index type.
Their simultaneous satisfaction establishes that the hypotheses are consistent,
and identifies no physical mechanism common to a double well, a register, a mesh
and a cortical cover. In particular nothing here interprets \texttt{Cortex} as
cortex.

\subsection{Four active branches, not one}
The active chain replaces the passive predictive node with a heat-budget node.
The development discharges its hypotheses four times over, at increasing
physical specificity: with a noisy actuator, with that actuator held to a named
register budget, with a process-dependent actuated field, and with a kernel
built from mode occupancies on a product space.
\formal{chain-active-hypotheses-jointly-satisfiable}
\formal{chain-active-budget-jointly-satisfiable}
\formal{chain-active-actuated-jointly-satisfiable}
\formal{chain-active-microscopic-jointly-satisfiable}
\scope Only the arrow feeding coarse-graining differs between them; every
downstream input is unchanged. That four arrangements discharge the same
obligations is evidence that the obligations are weak, not that the arrangements
agree about a mechanism.

\subsection{The obligations are not free}
Two kinds of rejection fence the active node. The first concerns the budget.

\begin{claim}
The actuator does not satisfy \texttt{ActiveBound} at budget $0$.
\end{claim}
\begin{proof}
\texttt{ActiveBound} requires the process's cost to lie below the budget. The
actuator's cost is positive, and a positive real is not at most zero.
\end{proof}
\formal{thermalagency-zero-budget-rejected}

The second concerns the limit. Each convergence arrow says that, \emph{given}
the bound, a named energy sequence tends to a named target. Since the bound
genuinely holds for these processes, the arrow can be applied, and the target is
then pinned by uniqueness of limits.

\begin{claim}
The active convergence arrow fails for the constant-zero energy sequence at
target $3$; for the actuated sequence at target $0$; and for the microscopic
sequence at target $0$. The passive arrow likewise fails for the remembered
sequence at target $0$.
\end{claim}
\begin{proof}
In each case supply the arrow with the bound that does hold, obtaining
convergence of the sequence to the stated target. The same sequence is already
known to converge --- the constant sequence to $0$, the others to their computed
spatial limits. Uniqueness of limits in $\mathbb{R}$ equates the two values, and
the resulting numerical identity is false.
\end{proof}
\formal{thermalagency-wrong-limit-rejected}
\formal{actuated-wrong-limit-rejected}
\formal{microscopic-wrong-limit-rejected}
\formal{microscopic-passive-wrong-limit-rejected}
\scope These are rejections rather than vacuous implications precisely because
the premise is inhabited in each case. They show that the convergence arrows
constrain which energy sequence and which limit may be supplied; they do not
show that any particular sequence is the physically correct one. Choosing it
remains a modelling obligation outside the theorem.

\subsection{The passive edge does not pick out its hardware}
\begin{claim}
Two different declared systems over the same joint law --- one retaining its
memory, one scrambling it --- discharge the identical passive convergence edge.
\end{claim}
\begin{proof}
Each supplies the joint law that the edge's construction requires, and the
convergence itself is the same spatial fact about the shared energy sequence.
Neither proof consults the system's efficiency.
\end{proof}
\formal{microscopic-e45}
\formal{microscopic-e45-scrambled}
\scope The edge is therefore insensitive to how well the hardware preserves
what it stores. No configuration law, partition or efficiency is selected by
discharging it, which bounds how much the passage through this node can be read
as a claim about physical implementation.
